2. Saddha - Confidence & Lucidity

The Pali word 'saddha' has more than one meaning in Buddhism. In the scriptures, faith or saddha is professing the confidence in and the sense of assurance based on understanding that one places on the Buddha, Dhamma, and Sangha. It is not the blind belief based on wrong view or ignorance. IT is compared to the gem of the Universal Monarch, which when put into a bowl of muddy water immediately causes the mud to settle to the bottom and the water becomes sparkling and drinkable. Likewise when saddha arises in the mind, the impure mental states disappear and the mind becomes clear or lucid, purified like sparkling drinkable water. Due to mental purification by saddha, one will not act greedily, or be filled with anger and hatred, or believe in wrong views and superstitious beliefs but place one's trust in the teachings of the Perfectly Englightened One.

Therefor the word 'saddha' has two qualities, namely: confidence and lucidity.

The saddha of the Lotus-Like Lay Buddhist consists of the following:
1. Confidence in the Buddha (Perfectly Englightened One)
2. Confidence in the Dhamma (Teachings of the Buddha)
3. Confidence in the Sangha (Order of Bikkhus/Bikkhunis)
4. Confidence in Kamma (Cause of Volitional Action)
5. Confidence in Vipaka (Results of Kamma)

Confidence in Kamma and Vipaka is the consequence of one's confidence in the Buddha, Dhamma and Sangha of the Three Gems.

The Commentary mentions four kinds of confidence (saddha):

2.1 Confidence in One's Own Destiny (Agama-saddha)

This is the confidence of one who is bent on attaining Englightenment and becoming a Perfectly Englightened Buddha. The excellence of the Bodhisatta is unlimited. He has firm determination and cannot be shaken. Ledi Sayadaw says that after receiving sure prediction of becoming a Supreme Buddha, the Bodhisattahas unwavering faith in the Triple Gem. This means that he never doubts that merit comes through doing good deads by one's own effort.

2.2 Confidence through Attainment (Adhigama-saddha)

This is the confidence of Noble individuals, i.e., the Ariyas. Their faith is unshakeable from the moment they realized Sainthood and will remain so in whatever existences they are reborn.

In 'A Manual of the Excellent Man'[1], Ven. Ledi Sayadaw wrote: "A stream-winner (Sotapanna) may be born into a non-Buddhist family but will not be led into professing another religion even on pain of instant death. They would rather die than forsake their firm confidence in the Buddha's teaching. This confidence never falters, but grows until they attain Nibbana."

There is another passage that describes a Buddhist: "One is a satisfactory Buddhist, if one becomes indignant at being called an adherent of another religion, and is pleased to be called a Buddhist." In other words, one is pleased to hear the Buddha's teachings extolled and displeased to hear another religion recommended.

In the Majjhima Nikaya Bahudhatuka Sutta, the Buddha has delcared: "Bhikkhus, it is impossible for one who has attained the Noble Right View to acknowledge another teacher as his/her teacher." This means that thos Buddhists who have attained Sainthood will not put the founders of all other religions on equal level with the Perfectly Englightened One, because there is no being who can compare with him in Virtue, Concentration and Wisdom except another Perfectly Englightened One!

Story of Sura Ambattha, chief among male lay disciples of Unshakeable Faith

In Anguttara Nikaya Commentary, the layman Sura Ambattha had invited the Buddha to his house for a meal. After taking the meal, the Buddha delivered a Dhamma-talk upon which Sura Ambattha attained the state of Stream-winner (Sotapanna). After that, the Buddha departed. A moment later, he saw the Buddha at his doorstep again. It was Mara in the guide of the Buddha. Being somewhat perplexed at seeing the Buddha on his doorstep so soon, Ambattha enquired the reason for the second visit and was told: "When I told you about the Dhamma, I made one thoughtless mistake. I said that all five aggregates of the human being are impermanent, suffering and without self. But they are not all like that; some are permanent, pleasurable and with self."

Sura Ambattha thought: "Buddhas will never say anything thoughtlessly. This person is definitely not the Buddha. He must be Mara, the antithesis of the Buddha." So Sura Ambattha answered: "I can see through your disguise. You are Mara."

When Mara heard this, he knew that he had been exposed and admitted: "Yes, I am Mara."

Sura Ambattha answered: "Even if a thousand Maras like you came here, my faith cannot be shaken. I have realized the Dhamma taught by the Blessed One that all formations are impermanent, suffering and non-self. Now, depart from my doorstep!"

When Mara heard these words, he was unable to reply; defeated, he left immediately. Due to Sura Ambattha's confidence, the Buddha singled him as pre-eminent (etadagga) in faith among the male lay disciples.

2.3 Established Confidence (Okappana-saddha)

This is the confidence after having settled on and having put one's trust in the Three Gems. It is defined as faith, which is inspired by the Noble Attributes of the Three Gems. It is firm and lasts for one's whole life. After death, however, it vanishes from one's consciousness. To enrich one's faith requires understanding the Noble Attributes of the Triple Gem. The recommended method of development of saddha is to recollect the Virtues of the Buddha, the Dhamma and the Sangha and through the practice of Satipatthana Vipassana meditation to attain Right View.

2.4 Confidence through Belief (Pasada-saddha)

This is revering the appearance and believing. It is confidence in the Three Gems because the Buddha, Dhamma, and Sangha are recognized as being worthy of reverence. It is based upon a superficial high regard for the Three Gems and not based on a deep conviction. There are many people born in Buddhist countries, who do not attain even this first step. According to the Dhammapada Verse 2, just before Matthakundali died, the Buddha appeared before him out of compassion. Seeing the Buddha, he was pleased, developed faith and dying with a pure heart, was reborn in Heaven.

Story of Matthakundali, who attained to the Deva realm through Faith[2]

While residing at the Jetavana monastery in Savatthi, the Buddha uttered Vers (2) of the Dhammapada, with reference to Matthakundali, a young Brahmin. Matthakundali was a young Brahmin whose father, Adinnapubbaka, was very stingy and never gave anything in charity. Even the gold ornaments for his only sone were made by himself to save payment for workmanship. When his son fell ill, no physician was consulted, until it was too late. When he realized that his son was dying, he had the youth carried outside on to the verandah, so that people coming to his house would not see his possessions. On that morning, the Buddha arising early from his deep meditation of compassion saw, in his Net of Knowledge, Matthakundali lying on the verandah. So when entering Savatthi for almsfood with his disciples, the Buddha stood near the door of the Brahmin Adinnapubbaka. The Buddha sent forth a ray of light to attract the attention of the youth, who was facing the interior of the house. The youth saw the Buddha; and as he was very weak he could only profess his faith mentally. But that was enough. When he passed away with his heart in devotion to the Buddha he was reborn in the Tavatimsa heaven.

From his celestial abode young Matthakundali, seeing his father mourning over him at the cemetery, appeared to the old man in the likeness of his old self. He told his father about his rebirth in the Tavatimsa world and also urged him to approach and invite the Buddha to a meal. At the house of Adinnapubbaka, the question of whether one could or could not be reborn in a celestial world simply by mentally professing profound faith in the Buddha, without giving in charity or observing the moral precepts, was brought up. So the Buddha willed that Matthakundali should appear in person; Matthakundali soon appeared fully decked with celestial ornaments and told them about his rebirth in the Tavatimsa world. Then only, the audience became convinced that the son of the Brahmin Adinnapubbaka by simply devoting his mind to the Buddha had attained much glory.

2.5 Devotion in Buddhism

Devotion or adoration (Pali 'bhatti'; Sanskrit 'bhakti') through worship of objects or symbols connected with the founder is common in all religions. The Buddha repeatedly discouraged any excessive blind veneration paid to him personally because he know that too much emotional devotion would obstruct mental development in the practice of the Noble Eightfold Path.

In the Scriptures, there is a story of the monk Vakkali who was full of devotion and love for the Buddha. Even when he was gravely ill, all he wanted was to see the Buddha, as he was ever so desirous to behold Him bodily, To him the Buddha said: "What good will it be to see this foul body? Vakkali, he who see the Dhamma (Teaching) sees me. Indeed, Vakkali, seeing the Dhamma is seeing me, seeing me is seeing the Dhamma." (Khanda Samyutta, Vakkali Sutta).

In another incident shortly before His decease (Parinibbhana), the twin Sala trees broke out in full bloom and their blossoms rained upon His body together with celestial flowers and scented powder from the sky while the sound of heavenly music and voices filled the air out of reverence for Him, the Buddha advised:

"Yet not thus, Ananda is the Tathagata respected, venerated, esteemed, worshipped and honoured in the highest degree. But, Ananda, whatsoever bhikkhu or bhikkhini, layman or laywoman abides by the Teaching, lives uprightly in the Teaching, walks in the way of the Teaching, it is by him that the Tathagata is respected, venerated, esteemed, worshipped and honoured in the highest degree." (Mahaparinibbhana Sutta, V, 6).

These two teachings appear to convey the message of the Master that a true and deep understanding of the Dhamma through practice of the Noble Eightfold Path is vastly superior to any external homage or mere emotional devotion. However, we should not conclude that the Buddha disparaged a reverential and devotional attitude of mind stemming from a true understanding and deep admiration of what is great and noble. It should be emphasized that 'Seeing the Dhamma' is not mere intellectual knowledge through study or logic but the experiential knowledge (bhavanamaya nana) through the practice of the Noble Eightfold PAth. This is the saddha that affirms that no system wherein the Noble Eightfold Path is absent can claim to lead to Emancipation.

The Buddha refused to be made an object of a personality cult. (Note: No Buddha images appeared during his lifetime and the first five centuries after Parinibbana). But he knew that respect and homage paid to those worthy of it is a great blessing, which he preached in the Mangala Sutta. This is because when one assumes a respectful attitude, one develops humility and is able to recognize the superior qualities of others and learn from them. This is on of the factors necessary for progress, whether spiritual or worldly.

Devotion is a natural expression of confidence (saddha) and plays a vital role in the balanace of the facultires (indriya) with its complement faculty of wisdom (panna). A one-sided development of the intellectual faculties (intelligence, insight, wisdom) tends to make one skeptical while a one-sided development of faith or devotion tends to make one gullible. Hence both faculties must be balances in order to lead to mental and spiritual progress.

2.6 Benefits of Recollection of the Buddha, Dhamma, and Sangha

For the laity, the routine acts of devotion like the offering of flowers, lights, incense and worship (puja) are a means to connect with the Three Gems in their daily life. More important and of greater validity than these outward forms of devotion is the capacity to turn them to the practice of devotional mental culture such as the contemplation of the Virtues of the Buddha that will lead to great benefits, namely:

i) Acquires abundant faith, which purifies the mind so that mindfullness and concentration are easily established

ii) Productive of joy, which is helpful in difficult times e.g. sickness, loss or facing hardships

iii) Instills confidence in oneself thereby dispelling fear, anxiety, doubt and restlessness.

Devotional meditation such as the Recollection of the Buddha, Dhamma, and Sangha can serve as an invaluable aid in attaining concentration, which is the basis of liberating insight in the practice of Satipatthana Vipassana meditation.